% Canonical TeX edition. Initial text inherited from 3b067cc2eba4ac2f2f91248448727639d8ac7c5a
% Source: research/case_study/phase277_interval_covers.md
% Subsequent mathematical edits require a new audit.
% Options for packages loaded elsewhere
\PassOptionsToPackage{unicode}{hyperref}
\PassOptionsToPackage{hyphens}{url}
\documentclass[
  11pt,
  a4paper,
]{article}
\usepackage{xcolor}
\usepackage[margin=1in]{geometry}
\usepackage{amsmath,amssymb}
\setcounter{secnumdepth}{-\maxdimen} % remove section numbering
\usepackage{iftex}
\ifPDFTeX
  \usepackage[T1]{fontenc}
  \usepackage[utf8]{inputenc}
  \usepackage{textcomp} % provide euro and other symbols
\else % if luatex or xetex
  \usepackage{unicode-math} % this also loads fontspec
  \defaultfontfeatures{Scale=MatchLowercase}
  \defaultfontfeatures[\rmfamily]{Ligatures=TeX,Scale=1}
\fi
\usepackage{lmodern}
\ifPDFTeX\else
  % xetex/luatex font selection
  \setmainfont[]{STIXGeneral}
  \setmonofont[]{Menlo}
\fi
% Use upquote if available, for straight quotes in verbatim environments
\IfFileExists{upquote.sty}{\usepackage{upquote}}{}
\IfFileExists{microtype.sty}{% use microtype if available
  \usepackage[]{microtype}
  \UseMicrotypeSet[protrusion]{basicmath} % disable protrusion for tt fonts
}{}
\makeatletter
\@ifundefined{KOMAClassName}{% if non-KOMA class
  \IfFileExists{parskip.sty}{%
    \usepackage{parskip}
  }{% else
    \setlength{\parindent}{0pt}
    \setlength{\parskip}{6pt plus 2pt minus 1pt}}
}{% if KOMA class
  \KOMAoptions{parskip=half}}
\makeatother
\setlength{\emergencystretch}{3em} % prevent overfull lines
\providecommand{\tightlist}{%
  \setlength{\itemsep}{0pt}\setlength{\parskip}{0pt}}
\usepackage{newunicodechar}
\newfontfamily\numeralfallback{STIX Two Text}
\newunicodechar{∎}{\ensuremath{\blacksquare}}
\usepackage{fvextra}
\DefineVerbatimEnvironment{Highlighting}{Verbatim}{breaklines,breakanywhere,breakanywheresymbolpre={},breakanywheresymbolpost={},breaksymbolleft={},breaksymbolright={},commandchars=\\\{\}}
\newunicodechar{⁰}{{\numeralfallback\char"2070\relax}}
\newunicodechar{₀}{{\numeralfallback\char"2080\relax}}
\newunicodechar{₁}{{\numeralfallback\char"2081\relax}}
\newunicodechar{₂}{{\numeralfallback\char"2082\relax}}
\newunicodechar{₃}{{\numeralfallback\char"2083\relax}}
\newunicodechar{₄}{{\numeralfallback\char"2084\relax}}

\usepackage{seqsplit}
\renewcommand{\texttt}[1]{{\ttfamily\seqsplit{#1}}}
\setlength{\emergencystretch}{6em}

\newunicodechar{⁻}{{\numeralfallback\char"207B\relax}}
\newunicodechar{⁸}{{\numeralfallback\char"2078\relax}}
\newunicodechar{ⁱ}{{\numeralfallback\char"2071\relax}}
\newunicodechar{⁴}{{\numeralfallback\char"2074\relax}}
\newunicodechar{⁷}{{\numeralfallback\char"2077\relax}}
\newunicodechar{⁶}{{\numeralfallback\char"2076\relax}}
\newunicodechar{⁵}{{\numeralfallback\char"2075\relax}}
\newunicodechar{₈}{{\numeralfallback\char"2088\relax}}
\newunicodechar{ᵢ}{\textsubscript{i}}
\newunicodechar{ⱼ}{\textsubscript{j}}
\newunicodechar{ₜ}{\textsubscript{t}}
\usepackage{bookmark}
\IfFileExists{xurl.sty}{\usepackage{xurl}}{} % add URL line breaks if available
\urlstyle{same}
% fallback for those not using the hyperref driver hyperxmp:
\makeatletter
\@ifundefined{xmpquote}{\newcommand{\xmpquote}[1]{#1}}{}
\makeatother
\hypersetup{
  hidelinks,
  pdfcreator={LaTeX via pandoc}}

\author{}
\date{}

\begin{document}

\section{Lifting minimal covers of an
interval}\label{lifting-minimal-covers-of-an-interval}

Private continuation, phase 277, 11 October 2026. The rank-four theorem
used below is the accepted theorem of phase 256, independently reviewed
in phase 273 and published in phase 274. Nothing here proves
finite-cover multiplicity when the only intersection defect has rank at
least five.

\subsection{1. The extension supplied by the finite-cover
construction}\label{the-extension-supplied-by-the-finite-cover-construction}

Let \(G=\langle g_0,\ldots,g_{n-1}\rangle\) be a finite string group.
All maps in this note preserve the ordered marks. For a contiguous
interval \(I=[i,j]\), write \(G_I=\langle g_i,\ldots,g_j\rangle\).

\textbf{Lemma 1.} Let \(Q\to G_{[i,j]}\) be a finite string C-group
cover. There is a finite string C-group cover \(R\to G\) whose marked
subgroup \(R_{[i,j]}\to G_{[i,j]}\) is isomorphic to the given cover.

\textbf{Proof.} We spell out the preservation assertion in the
construction of Monson--Schulte, Lemma 3.3, rather than use only its
finite-existence conclusion. Suppose first that \(i=0\) and \(j<n-1\). A
finite regular extension \(E\) of \(Q\), supplied by their Lemma 3.1,
has an initial maximal parabolic identified with \(Q\). Form the
subgroup \[
 R'=\langle(g_t,e_t):0\le t\le j+1\rangle
       \le G_{[0,j+1]}\times E.
\] The initial subgroup generated by the first \(j+1\) marks is the
graph of the epimorphism \(Q\to G_{[0,j]}\), and is therefore exactly a
copy of \(Q\). Projection to \(E\) is an epimorphism and is injective on
that maximal parabolic. The standard quotient criterion proves that
\(R'\) is a string C-group. Projection to the first factor covers
\(G_{[0,j+1]}\), and both factors are finite. This is the actual proof
of Monson--Schulte's lemma. In particular it preserves the given cover
\(Q\), not merely those smaller parabolics on which its map happens to
be injective.

Starting from an arbitrary interval \([i,j]\), apply the same
construction to the consecutively relabelled string
\(g_i,\ldots,g_{j+1}\). Repeat to extend the right endpoint up to
\(n-1\). Then reverse the order of the marks and apply the construction
repeatedly to extend the left endpoint down to zero. Each extension
preserves the previously constructed maximal parabolic exactly, and
therefore preserves its subgroup on the original interval. The composite
epimorphisms preserve all marks. ∎

The classical input is the seven-page author version of Monson--Schulte:
Lemma 3.1 on PDF pages 4--5, and the direct-product construction in the
proof of Lemma 3.3 on PDF page 5. Its Theorem 3.4 gives existence but,
by itself, would conceal the preservation assertion needed here. The
source bytes are bound in phase277\_sources.json. This is a use of prior
mathematics, not a new finite-cover existence claim.

\subsection{2. Minimal descendants retain a minimal
interval}\label{minimal-descendants-retain-a-minimal-interval}

A finite string C-group cover \(Q\to L\) is minimal if every string
C-group quotient still covering \(L\) is the same marked cover.

\textbf{Theorem 2.} Every finite minimal string C-group cover of a
contiguous parabolic \(G_I\) occurs as the \(I\)-parabolic of some
finite minimal string C-group cover of \(G\).

\textbf{Proof.} Begin with a minimal cover \(Q\to G_I\) and the
extension \(R\to G\) from Lemma 1. Among the string C-group quotients of
\(R\) still mapping onto \(G\), choose a minimal one, denoted \(M\).
This is possible because \(R\) is finite. Every further C-group quotient
of \(M\) over \(G\) would be a quotient of \(R\), so \(M\) is minimal
among all C-covers of \(G\), not just within a smaller list.

The quotient map restricts to a surjective marked homomorphism \[
 Q\cong R_I\longrightarrow M_I\longrightarrow G_I.
\] A standard parabolic of a string C-group is a string C-group. Thus
\(M_I\) is a C-group quotient of \(Q\) still covering \(G_I\).
Minimality of \(Q\) forces the first map to be an isomorphism. ∎

\textbf{Corollary 3.} If \(G_I\) has at least \(m\) pairwise
nonisomorphic finite minimal C-covers, then \(G\) has at least \(m\),
for every positive integer \(m\). If \(G_I\) has infinitely many, so
does \(G\).

\textbf{Proof.} Apply Theorem 2 to each local cover. A marked
isomorphism of two resulting global groups restricts to a marked
isomorphism of their \(I\)-parabolics. Distinct local covers therefore
produce distinct global ones. The same argument works for every finite
subfamily of an infinite family. ∎

This does not assert that the parabolic of an arbitrary global minimal
cover must be minimal. The theorem constructs global minimal covers
whose specified parabolic is the prescribed local minimal one.

\subsection{3. A consequence of the rank-four theorem in every
rank}\label{a-consequence-of-the-rank-four-theorem-in-every-rank}

\textbf{Corollary 4.} Let \(P\) be any finite polytope, of arbitrary
rank, and let \(G=\operatorname{Mon}(P)\). If some subgroup generated by
four consecutive marks is not a string C-group, then \(P\) has
infinitely many finite minimal regular polytopal covers.

\textbf{Proof.} Let \(I=[i,i+3]\) be such an interval. Each endpoint
mark lies outside the group generated by the other three: it lies
outside the group generated by all other marks of \(G\), by the
flag-face argument. The four marks are distinct nonidentity involutions
and satisfy the required string relations. Hence Theorem 1 of phase 256
applies directly to the finite marked group \(G_I\); no separate
assertion that \(G_I\) is itself a polytope's monodromy group is needed.

For every prime \(p>4|G_I|\), that theorem gives a finite C-cover with
\(p\)-group kernel. A minimal C-group quotient still covering \(G_I\)
has nontrivial \(p\)-power kernel, since \(G_I\) is not C. These local
minimal covers have distinct orders for distinct primes. Corollary 3 now
gives infinitely many global finite minimal C-covers of \(G\). The usual
normal flag-kernel correspondence turns them into finite minimal regular
polytopal covers of \(P\). ∎

No bound or prime-power assertion about the \emph{global} cover degree
is made. The extension procedure can introduce other primes.
Distinctness is read from the indicated interval parabolic, not inferred
from the total orders.

More generally, any later finite-multiplicity theorem for a marked
interval can be transported by Theorem 2 without repeating this
argument.

\subsection{4. The genuinely unresolved
case}\label{the-genuinely-unresolved-case}

Independent rank-three string groups need not be C-groups. The familiar
rank-three monodromy theorem is a statement about actual polytopes and
cannot be applied to this broader class. There is a direct reduction to
phase 256 which supplies the needed rank-three result instead. For a
finite independent rank-three group \(L=\langle a,b,c\rangle\), append
an independent central involution \(t\), obtaining \(L\times C_2\) with
marks \((a,1),(b,1),(c,1),(1,t)\). Both extreme independence conditions
of phase 256 hold. Restrict each large-prime \(p\)-kernel C-cover
supplied by phase 256 to its first three marks. This parabolic is a
C-group mapping onto \(L\), with a \(p\)-group kernel. If \(L\) is
non-C, its finite minimal descendants have nontrivial \(p\)-power
degrees, and distinct primes give infinitely many. Ranks one and two
already satisfy the intersection property.

Now suppose a finite non-C string group with independent marks has only
finitely many finite minimal C-covers, and choose such an example of
least rank. All its proper contiguous parabolics must be C-groups.
Otherwise a proper non-C interval has infinitely many finite minimal
covers by minimality of the rank, and Corollary 3 contradicts finiteness
for the whole group. Thus the smallest unresolved rank for this stronger
infinitude statement is at least five, and the obstruction may be
assumed to be the top intersection \[
 G_{[0,n-2]}\cap G_{[1,n-1]}\ne G_{[1,n-2]},
\] with both end groups already C. This is a reduction in the broader
class of independent marked groups; it should not silently identify
every such group with an actual polytope monodromy group.

The analogous conclusion from a least-rank counterexample restricted
only to actual polytopes needs care: a parabolic of their monodromy
group acts on the disjoint union of many sections and need not have been
separately identified as the monodromy group of one polytope. Corollary
4 avoids that issue because phase 256 is a theorem for the broader
marked-group class.

The locally faithful amalgam criterion in the previously reviewed phase
246 manuscript remains relevant to this remaining case. It does not by
itself suffice: the Tomotope already shows that locally faithful
odd-primary covers may not exist. Rank four was resolved by enlarging
the end groups before the second homology cover. A higher-dimensional
substitute for that end enlargement, with simultaneous control of all
required intersections, is the outstanding mathematical step.

\end{document}
